\section{K-011: all two-primary lifts in dimension five}

\begin{theorem}[K-011]
For every nonzero polystable tensor in $(\mathbb C^5)^{\otimes3}$,
the two-primary part of its determinant stabilizer has order at most
$16$.  This bound is attained.  Every grading-fixed quotient point
for a primitive sixteenth root has, after diagonalization and scalar
normalization, one of the seventy-seven spectral types in K-010.
No additional types arise from order-$32$ or order-$64$ lifts.
\end{theorem}

\begin{proof}
K-009 gives a two-primary lift cutoff of $64$ in dimension five.
The primitive-character calculation K-010 excludes order-$32$
primitive lifts and, by reduction, order-$64$ primitive lifts.
To settle nonfaithful lifts, the prescribed-character version of the
refinement lemma is implemented in \texttt{refine\_character.cpp}
and \texttt{recursive\_character.py}.  For target $q$, every stage
modulo $M$ retains all normalizing odd units satisfying
$u d(W)=q\pmod M$, including every odd unit when the target is zero.
The target is reduced at every coarse stage; it is never replaced
by a primitive target.  The exact numbers of canonical candidates
and feasible triples are as follows:
\[
\begin{array}{c|rr|rr}
 &\multicolumn{2}{c|}{q=2}&\multicolumn{2}{c}{q=4}\\
 M&\text{candidates}&\text{feasible}&\text{candidates}&\text{feasible}\\\hline
 2&8&6&8&6\\
 4&154&40&162&109\\
 8&9051&506&10327&1047\\
 16&261715&718&395519&3298\\
 32&392496&77&1396339&1266\\
 64&36878&0&616028&77
\end{array}                                                   \tag{K-011a}
\]
Every feasibility decision is accompanied by an exact rational point
or an exact rational Farkas separator.  Floating-point linear algebra
is used only to propose certificates, whose rational equations and
inequalities are verified before acceptance.  The standalone
integer/\texttt{Fraction} checker is
\texttt{verify\_recursive\_certificates.py 5 2} or, respectively,
\texttt{verify\_recursive\_certificates.py 5 4}.  It also checks that
each certificate has exactly the corresponding spectrum in the
complete generated list.  There are $700302$ and $2418383$
certificates respectively.

A determinant character of order $32$ could have an order-$32$
primitive lift or an order-$64$ lift with target exponent $2$.
Both are excluded.  Order $64$ is excluded by K-010, and larger
two-power orders are excluded by K-009.  Order $16$ is attained by
the rechecked chain tensor P-F60, or by any representative of K-010.

For character order $16$, possible lift orders are $16,32,64$.
The seventy-seven feasible spectra at $(M,q)=(32,2)$ are all even;
division by two gives, in order, exactly the seventy-seven primitive
modulus-sixteen spectra.  The seventy-seven spectra at
$(M,q)=(64,4)$ are all divisible by four; division by four gives
the same list.  These equalities are checked directly in the
recorded integer spectra.  Odd powers and scalar shifts used in
normalization preserve the fixed tensor space.  Hence these larger
lifts contribute no new semistable fixed spaces.

Finally a nonzero point of the invariant quotient fixed by the
grading element has a closed tensor orbit and a product-linear
stabilizer lift with the prescribed determinant character, as in
K-009.  Its semisimple two-primary lift diagonalizes.  Torus
semistability supplies the necessary margins, so the complete
spectral classification applies.  K-010 proves that each of these
fixed spaces has one centralizer orbit of semistable tensors.
This proves the asserted coverage of all nonzero fixed quotient
points.
\end{proof}
